\documentclass{aastex631}

\usepackage{amsmath,amssymb}

\newcommand{\code}[1]{\texttt{#1}}
\newcommand{\Rem}{\mathrm{Re}_m}
\newcommand{\Lc}{L_c}
\newcommand{\Cp}{c_P}
\newcommand{\HT}{H_T}
\newcommand{\Mearth}{M_\oplus}
\newcommand{\Fconv}{F_{\rm conv}}

\shorttitle{Dynamo Criterion in ORCHARD}
\shortauthors{Tejada Arevalo}

\begin{document}

\title{A Magnetic Reynolds Number Diagnostic for Rocky Planets in ORCHARD}

\author{Roberto Tejada Arevalo}

\begin{abstract}
We document a lightweight post-processing diagnostic that evaluates whether, where,
and for how long a dynamo can operate in the interior of an ORCHARD rocky-planet
model.  We follow the method of the CMAPPER code \citep[][Section~3.3]{Zhang2025}:
in every liquid cell we form the magnetic Reynolds number $\Rem=\mu_0 v \Lc \sigma$
from quantities ORCHARD already produces (convective luminosity, density, melt
fraction, and equation-of-state thermodynamics) and flag the cell as dynamo-capable
where $\Rem$ exceeds a critical value $\approx50$.  Applied to a $1\,\Mearth$
super-Earth, the liquid iron core reaches $\Rem\sim10^4$ --- consistent with the
Earth's core --- and the calculation reproduces Figure~2 of \citet{Zhang2025}.  This
note is intended as a recipe a student can implement; the reference implementation is
\code{magnetic\_reynolds.py}.
\end{abstract}

\section{The dynamo criterion}

A planetary dynamo requires a region of convecting, electrically conducting fluid that
stirs the magnetic field faster than ohmic diffusion erases it.  That ratio is the
magnetic Reynolds number,
\begin{equation}
  \Rem = \mu_0\, v\, \Lc\, \sigma ,
  \qquad\text{dynamo active if}\quad \Rem > \Rem^{\rm crit}\approx 50 ,
  \label{eq:rem}
\end{equation}
where $\mu_0=4\pi\times10^{-7}\,\mathrm{H\,m^{-1}}$, $v$ is the convective flow speed,
$\Lc$ is the thickness of the liquid (convecting) region, and $\sigma$ is its electrical
conductivity.  The critical value $\approx 50$ is the self-sustaining threshold of
\citet{ChristensenAubert2006}.  Equation~\eqref{eq:rem} is dimensionless; we assemble
its ingredients in CGS (ORCHARD's native units) and convert $v$ and $\Lc$ to SI at the
end, since $\mu_0$ and $\sigma$ are defined in SI.

\section{The three ingredients}

\paragraph{Length $\Lc$.}  ORCHARD outputs a melt fraction per cell; we call a cell
``liquid'' when \code{melt\_fraction}~$\geq 0.75$ (a provisional threshold).  $\Lc$ is
the radial extent of a contiguous block of liquid cells, taken from the cell-boundary
radii \code{r\_b}.  The mantle (magma ocean) and the iron core are treated as separate
blocks.  Typical value: $\Lc\sim3500~\mathrm{km}=3.5\times10^{8}~\mathrm{cm}$.

\paragraph{Speed $v$.}  We use the mixing-length scaling of \citet{Christensen2010},
adopted by CMAPPER,
\begin{equation}
  v = \left(\frac{\Lc\,\Fconv}{\rho\,\HT}\right)^{1/3},
  \qquad
  \HT = \frac{\Cp}{\alpha\,g} ,
  \label{eq:v}
\end{equation}
where $\HT$ is the temperature scale height (along an adiabat
$\mathrm{d}T/\mathrm{d}r|_s=-\alpha g T/\Cp$, so $\HT=T/|\mathrm{d}T/\mathrm{d}r|$).
The local gravity is $g=GM_{\rm enc}/r^2$, with the enclosed mass $M_{\rm enc}$ from
\code{m\_b} and the radius from \code{r}.  The specific heat $\Cp$ and thermal
expansivity $\alpha$ come from the relevant equation of state (Mg$_2$SiO$_4$ for the
mantle, the Gonz\'alez iron tables for the core).  The convective \emph{flux} $\Fconv$
is discussed in Section~\ref{sec:units}.

\paragraph{Conductivity $\sigma$.}  ORCHARD models \emph{thermal} conductivity, but a
dynamo cares about \emph{electrical} conductivity, so we bridge the two with the
Wiedemann--Franz law, $k=L_0\sigma T$.  For the magma ocean we invert the same combined
mantle conductivity ORCHARD uses internally --- the electronic term of
\citet{Stamenkovic2011} plus the lattice term of \citet{PengDeng2025} --- and weight by
the melt fraction:
\begin{equation}
  \sigma_{\rm sil} = x_{\rm melt}\,
  \frac{k_{\rm el}(T) + k_{\rm lat}(\rho,T)}{L_0\,T} ,
  \qquad
  \sigma_{\rm Fe} = \frac{k_c}{L_0\,T} ,
  \label{eq:sigma}
\end{equation}
with $L_0\approx2.5\times10^{-8}\,\mathrm{W\,\Omega\,K^{-2}}$ the Lorenz number.  For
the metallic iron core (right-hand expression, CMAPPER Eq.~18) we use the core thermal
conductivity $k_c$ directly; there Wiedemann--Franz is rigorous because electrons carry
both heat and charge.  Typical values: $\sigma_{\rm sil}\sim2\times10^4~\mathrm{S\,m^{-1}}$,
$\sigma_{\rm Fe}\sim10^5~\mathrm{S\,m^{-1}}$.

\section{Practical notes (the unit traps)}
\label{sec:units}

Two conversions decide the answer by many orders of magnitude:
\begin{enumerate}
\item \textbf{\code{conv\_flux} is a luminosity, not a flux.}  ORCHARD stores the
  convective luminosity ($\sim10^{21}~\mathrm{erg\,s^{-1}}$, comparable to the planet's
  total internal luminosity).  The per-area flux that enters Equation~\eqref{eq:v} is
  $\Fconv = \code{conv\_flux}/(4\pi r_b^2)\sim10^{3}~\mathrm{erg\,cm^{-2}\,s^{-1}}
  \;(\sim0.1~\mathrm{W\,m^{-2}})$.  Omitting the $4\pi r_b^2$ gives an unphysical
  $\sim10^{18}~\mathrm{W\,m^{-2}}$.
\item \textbf{The EOS expects pressure in GPa.}  ORCHARD's pressure is in
  $\mathrm{dyn\,cm^{-2}}$; multiply by \code{dyn\_to\_GPa}~$=10^{-10}$ before calling
  \code{get\_cp\_pt}/\code{get\_alpha\_pt}.  The returned $\Cp$ and $\alpha$ are already
  CGS.
\end{enumerate}
A cheap self-check catches both: ORCHARD already saves \code{cp}, so compute $\Cp$ from
the EOS and compare.  In our runs the ratio is $1.00$ to within a percent, confirming
the units \emph{before} any $\Rem$ is trusted.

\paragraph{Order-of-magnitude check (iron core).}  With
$\Fconv\sim10^{3}~\mathrm{erg\,cm^{-2}\,s^{-1}}$, $\rho\sim10~\mathrm{g\,cm^{-3}}$,
$\Cp\sim10^{7}~\mathrm{erg\,g^{-1}\,K^{-1}}$, $\alpha\sim10^{-5}~\mathrm{K^{-1}}$, and
$g\sim10^{3}~\mathrm{cm\,s^{-2}}$, we get $\HT=\Cp/(\alpha g)\sim10^{9}~\mathrm{cm}$ and,
from Equation~\eqref{eq:v}, $v\sim\mathrm{few}~\mathrm{cm\,s^{-1}}\sim10^{-2}~\mathrm{m\,s^{-1}}$.
Then, in SI,
\begin{equation}
  \Rem \sim
  \underbrace{1.3\times10^{-6}}_{\mu_0}\cdot
  \underbrace{10^{-2}}_{v}\cdot
  \underbrace{3.5\times10^{6}}_{\Lc}\cdot
  \underbrace{2\times10^{5}}_{\sigma}
  \;\approx\; 10^{4} .
\end{equation}

\section{Example: a $1\,\Mearth$ super-Earth}

Figure~\ref{fig:rem} shows $\Rem(P)$ for the thin-envelope $1\,\Mearth$ model at several
ages.  The structure matches Figure~2 of \citet{Zhang2025}: a magma-ocean branch
($P\lesssim115~\mathrm{GPa}$, $\Rem\sim10^{3\text{--}4}$) and an iron-core branch
($P\gtrsim115~\mathrm{GPa}$, $\Rem\sim10^{4}$), separated by a sharp jump at the
core--mantle boundary that is \emph{entirely} a conductivity contrast ($\sigma$ rises
from $\sim2\times10^4$ to $\sim10^5~\mathrm{S\,m^{-1}}$).  The iron-core $\Rem$ sits near
$10^4$, consistent with the Earth's core.  Following the cooling history, the iron core
is molten and convecting (and thus dynamo-capable) from $\sim1~\mathrm{Myr}$ to
$\sim2\text{--}3~\mathrm{Gyr}$, after which it solidifies and the core dynamo shuts off
(the $4.56~\mathrm{Gyr}$ curve has no iron-core branch); a shrinking magma-ocean dynamo
persists somewhat longer.

\begin{figure}[t]
\centering
\includegraphics[width=0.78\textwidth]{magnetic_reynolds_1Mearth.png}
\caption{Magnetic Reynolds number versus pressure for the ORCHARD $1\,\Mearth$
thin-envelope super-Earth at several ages, computed with the recipe of
Equations~\eqref{eq:rem}--\eqref{eq:sigma}.  Thin lines show the full liquid profile;
thick segments mark the liquid iron core.  Dashed: the dynamo threshold
$\Rem^{\rm crit}\approx50$; dotted: the Earth-core reference $\Rem\sim10^4$.  The
two-branch structure and the core value reproduce Figure~2 of \citet{Zhang2025}.}
\label{fig:rem}
\end{figure}

\paragraph{Caveat on the absolute value.}  The result is reliable to a factor of a few,
not better.  Mixing-length theory overestimates $v$ (rotation suppresses real core flows
toward $\sim$\,mm\,s$^{-1}$), while a modest $k_c$ underestimates $\sigma$; the errors
partly cancel near $10^4$.  This is inherent to the non-rotating scaling that CMAPPER
also adopts.  The robust outputs are therefore \emph{structural}: which layers exceed
$\Rem^{\rm crit}$, and the lifetime over which a region stays liquid and convecting.

\begin{thebibliography}{}
\bibitem[Christensen(2010)]{Christensen2010} Christensen, U.~R. 2010, \textit{Space Sci. Rev.}, 152, 565
\bibitem[Christensen \& Aubert(2006)]{ChristensenAubert2006} Christensen, U.~R., \& Aubert, J. 2006, \textit{Geophys. J. Int.}, 166, 97
\bibitem[Peng \& Deng(2025)]{PengDeng2025} Peng, Z., \& Deng, J. 2025, lattice thermal conductivity of liquid silicate
\bibitem[Stamenkovi\'c et al.(2011)]{Stamenkovic2011} Stamenkovi\'c, V., Breuer, D., \& Spohn, T. 2011, \textit{Icarus}, 216, 572
\bibitem[Zhang(2025)]{Zhang2025} Zhang, J. 2025, The CMAPPER Terrestrial Planet Evolution Code
\end{thebibliography}

\end{document}
